% !Mode:: "TeX:UTF-8"

\section{个人贡献}


\subsection{报告撰写}

本实验报告由本人独立完成，包括：
\begin{itemize}
    \item 系统设计原理阐述
    \item 四种边缘检测算法原理分析
    \item 修改内容详解
    \item 仿真结果整理与分析
    \item 资源利用率和时序结果分析
    \item 实验现象记录与总结
\end{itemize}


\subsection{编码工作}

\begin{itemize}
    \item 参与prewitt\_core.v模块的代码编写和调试
    \item 参与laplacian\_core.v模块的代码编写和调试
    \item 参与roberts\_core.v模块的代码编写和调试
    \item 协助修改hdmi\_bram\_sobel\_display.v的模式扩展
    \item 参与仿真测试用例的编写
\end{itemize}

\subsection{硬件连接}

\begin{itemize}
    \item 负责ZYNQ7020开发板与HDMI显示器的连接
    \item 负责USB串口线与PC的连接调试
    \item 负责上板验证和现象记录
\end{itemize}

\subsection{资本投入}

本次实验主要使用实验室提供的设备，个人无额外资本投入。


\subsection{开发问题描述与解决方案}

\subsubsection{Prewitt边缘幅度较低}

\begin{itemize}
    \item 问题描述：Prewitt模式下边缘图整体较暗，细节不如Sobel清晰
    \item 解决方法：降低阈值补偿，或理解这是算法特性（无2x权重）
\end{itemize}



\subsubsection{Laplacian噪声敏感}

\begin{itemize}
    \item 问题描述：Laplacian对噪声敏感，边缘图噪点较多
    \item 解决方法：提高阈值过滤噪声，或理解这是二阶导数的特性
\end{itemize}



\subsubsection{Roberts边缘断裂}

\begin{itemize}
    \item 问题描述：Roberts使用2×2卷积核，边缘有时不连续
    \item 解决方法：降低阈值保留更多边缘，或理解这是小卷积核的限制
\end{itemize}



\subsubsection{模式切换命令格式}

\begin{itemize}
    \item 问题描述：新增模式4/5/6后，需要确保上位机命令格式正确
    \item 解决方法：参考现有命令格式，扩展3-bit模式支持
\end{itemize}



\subsubsection{仿真测试覆盖}

\begin{itemize}
    \item 问题描述：需要确保新增模式的仿真测试完整性
    \item 解决方法：新增4项算法专项测试，覆盖边界条件
\end{itemize}



\subsubsection{资源利用优化}

\begin{itemize}
    \item 问题描述：新增3个算法模块可能增加资源消耗
    \item 解决方法：Prewitt/Roberts不消耗额外DSP，Laplacian可优化为左移
\end{itemize}
